% !TeX root = ../CV_Tramarin.tex

\sectioniten{Specifiche esperienze caratterizzate da attività di ricerca attinenti al settore concorsuale}{Other experiences characterized by research activities}

\subsectioniten{Partecipazione a Commissioni Valutazione di Dottorato}{Participation in PhD Evaluation Committees}

\cventryit{2025}{Commissione di valutazione per la Difesa pubblica della Tesi di Dottorato}{Dr. Daniel Bujosa Mateu}{\newline\emph{Supervisore}: Prof. Mikael Gidlund. \emph{Opponent}: Prof. Preben E. Mogensen, Aalborg University, Danimarca}{\newline Altri membri della commissione: Prof. Fortunato Santucci, University of L'Aquila, Prof. Erik Ström, Chalmers University of Technology, Sweden}{Mälardalen University, Västerås, Sweden}

\cventryen{2025}{Evaluation Committee Member for the Public Defense of the PhD Thesis}{Dr. Daniel Bujosa Mateu}{\newline\emph{Supervisor}: Prof. Mikael Gidlund. \emph{Opponent}: Prof. Preben E. Mogensen, Aalborg University, Denmark}{\newline Other committee members: Prof. Fortunato Santucci, University of L'Aquila, Prof. Erik Ström, Chalmers University of Technology, Sweden}{Mälardalen University, Västerås, Sweden}

\cventryit{2024}{Opponent nella Difesa pubblica della Tesi di Dottorato}{Dr. Pablo Gutiérrez Peón}{\newline\emph{Supervisore}: Prof. Elisabeth Uhlemann, Mälardalen University, Västerås, Svezia}{\newline Altri membri della commissione: Prof. Tingting Zhang, Mid Sweden University, Sweden; Prof. Soheil Samii, Linköping University, Sweden; Prof. Urban Bilstrup, Halmstad University, Sweden}{Mälardalen University, Västerås, Svezia}

\cventryen{2024}{Opponent for the Public Defense of the PhD Thesis}{Dr. Pablo Gutiérrez Peón}{\newline\emph{Supervisor}: Prof. Elisabeth Uhlemann, Mälardalen University, Västerås, Sweden}{\newline Other committee members: Prof. Tingting Zhang, Mid Sweden University, Sweden; Prof. Soheil Samii, Linköping University, Sweden; Prof. Urban Bilstrup, Halmstad University, Sweden}{Mälardalen University, Västerås, Sweden}

\cventryit{2023}{Presidente della Commissione esaminatrice per l'esame finale di Dottorato}{Corso di Dottorato di Ricerca in Ingegneria e Scienza dell’Informazione, Università di Siena}{Valutazione del candidato Dott. Giacomo Peruzzi}{19 Giugno 2023}{}

\cventryen{2023}{Chair of the Examination Committee for the Final PhD Examination}{PhD Program in Information Engineering and Science, University of Siena}{Evaluation of the candidate Dr. Giacomo Peruzzi}{June 19, 2023}{}

\cventryit{2021}{Commissione di valutazione per la Difesa pubblica della Tesi di Dottorato}{Dr. Luca Beltramelli}{\newline\emph{Supervisore}: Prof. Mikael Gidlund. \emph{Opponent}: Prof. Preben E. Mogensen, Aalborg University, Danimarca}{\newline Altri membri della commissione: Prof. Fortunato Santucci, University of L'Aquila, Prof. Erik Ström, Chalmers University of Technology, Sweden}{MidSweden University, Sundsvall, Svezia}

\cventryen{2021}{Evaluation Committee for the Public Defense of the PhD Thesis}{Dr. Luca Beltramelli}{\newline\emph{Supervisor}: Prof. Mikael Gidlund. \emph{Opponent}: Prof. Preben E. Mogensen, Aalborg University, Denmark}{\newline Other committee members: Prof. Fortunato Santucci, University of L'Aquila, Prof. Erik Ström, Chalmers University of Technology, Sweden}{MidSweden University, Sundsvall, Sweden}

\cventryit{2021}{Membro della Commissione esaminatrice per l'esame finale di Dottorato}{Corso di "Technology for Health", Università di Brescia}{Valutazione di 6 studenti di dottorato iscritti al corso}{26 Marzo 2021}{}

\cventryen{2021}{Member of the Examination Committee for the Final PhD Examination}{PhD Program in "Technology for Health", University of Brescia}{Evaluation of 6 PhD students enrolled in the program}{March 26, 2021}{}


\subsectioniten{Revisione di Progetti di Ricerca finanziati in ambito EU}{Expert Reviewer for the Evaluation of EU funded projects}

\cventryit{2024}{Revisore di Progetti di Ricerca Internazionali in ambito EURAMET}{European Association of National Metrology Institutes}{\newline nell'ambito delle Call ``Metrology Partnership'', Review Conference della sezione Digital Transformation (DIT), Bruxelles (Belgio)}{11-12 Novembre 2024}{Valutazione delle proposte progettuali ricevute nella call e selezione delle proposte da trasformare in progetti di ricerca europei finanziati in ambito metrologico}

\cventryen{2024}{Reviewer for International Research Projects within EURAMET}{European Association of National Metrology Institutes}{\newline within the framework of the ``Metrology Partnership'' Calls, Review Conference of the Digital Transformation (DIT) section, Brussels (Belgium)}{November 11-12, 2024}{Evaluation of the project proposals submitted to the call and selection of proposals to be developed into European-funded research projects in the field of metrology}

\cventryit{2023}{Revisore di Progetti di Ricerca Internazionali in ambito EURAMET}{European Association of National Metrology Institutes}{\newline  nell'ambito delle Call ``Metrology Partnership'', Review Conference della sezione Research Potential (RP), Amsterdam (Paesi Bassi)}{15-16 Novembre 2023}{Valutazione delle proposte progettuali ricevute nella call e selezione delle proposte da trasformare in progetti di ricerca europei finanziati  in ambito metrologico}

\cventryen{2023}{Reviewer for International Research Projects within EURAMET}{European Association of National Metrology Institutes}{\newline within the framework of the ``Metrology Partnership'' Calls, Review Conference of the Research Potential (RP) section, Amsterdam (Netherlands)}{November 15-16, 2023}{Evaluation of the project proposals submitted to the call and selection of proposals to be developed into European-funded research projects in the field of metrology}

\cventryit{2022}{Revisore di Progetti di Ricerca Internazionali in ambito EURAMET}{European Association of National Metrology Institutes}{\newline  nell'ambito delle Call ``Metrology Partnership'', Review Conference della sezione Health (HLT), Principato di Monaco}{14-15 Novembre 2022}{Valutazione delle proposte progettuali ricevute nella call e selezione delle proposte da trasformare in progetti di ricerca europei finanziati  in ambito metrologico}

\cventryen{2022}{Reviewer for International Research Projects within EURAMET}{European Association of National Metrology Institutes}{\newline within the framework of the ``Metrology Partnership'' Calls, Review Conference of the Health (HLT) section, Principality of Monaco}{November 14-15, 2022}{Evaluation of the project proposals submitted to the call and selection of proposals to be developed into European-funded research projects in the field of metrology}

\subsectioniten{Seminari su invito}{Invited Seminars}

\cventryit{2024}{Seminario su invito}
{``TSN in hybrid wired/wireless scenarios: measurement methods for latency characterization''}{Presentazione delle attività di ricerca relative al progetto PRIN 2022 PNRR ``Time Sensitive Networking for future enabling technologies:
metodi di misura e caratterizzazione metrologica in scenari ibridi
    cablati/wireless'' le cui attività sno nell'ambito del TC37 della IEEE IMS.}{24 Settembre 2024}{Mälardalen University, Västerås, Svezia}

\cventryen{2024}{Invited Seminar}
{``TSN in hybrid wired/wireless scenarios: measurement methods for latency characterization''}{Presentation of research activities related to the PRIN 2022 PNRR project ``Time Sensitive Networking for future enabling technologies: measurement methods and metrological characterization in hybrid wired/wireless scenarios,'' whose activities are within the scope of the IEEE IMS TC37.}{September 24, 2024}{Mälardalen University, Västerås, Sweden}

\cventryit{2021}{Seminario su invito}
{``Protocolli e standard IoT per la salute e il benessere''}{All'interno della giornata di studio ``Interfacce e Interazioni nel Medical Design''. Presentazione delle attività di ricerca relative al progetto FAR 2021 ``Smart Healthcare Modular System for Personalized Remote Health Monitoring'' le cui attività si inseriscono nell'ambito della sensoristica IoT per il biomedicale.}{19 novembre 2021}{Università Iuav di Venezia, Italia}

\cventryen{2021}{Invited Seminar}
{``Protocolli e standard IoT per la salute e il benessere''}{As part of the study day ``Interfacce e Interazioni nel Medical Design'' Presentation of research activities related to the FAR 2021 project ``Smart Healthcare Modular System for Personalized Remote Health Monitoring,'' whose activities fall within the scope of IoT sensor technology for biomedical applications.}{November 19, 2021}{Iuav University of Venice, Italy}

\cventryit{2016}{Seminario su invito}
{``High-Throughput Enhancements and Rate Adaptation Strategies for IEEE 802.11 Based Real-Time Communications''}{Presentazione delle attività di ricerca e discussione sulle tematiche di ricerca dell'analisi delle prestazioni di reti industriali di sensori wireless, svolte anche in seno alle linee di ricerca del TC-37 della IEEE IMS}{13 Giugno 2016}{Università di Verona, Dipartimento di Informatica}

\cventryen{2016}{Invited Seminar}
{``High-Throughput Enhancements and Rate Adaptation Strategies for IEEE 802.11 Based Real-Time Communications''}{Presentation of research activities and discussion on research topics related to the performance analysis of industrial wireless sensor networks, also conducted within the research lines of the IEEE IMS TC-37.}{June 13, 2016}{University of Verona, Department of Computer Science}


\sectioniten{Membro di Comitati Tecnici e Società Scientifiche}{Membership in International Societies and Technical Committees}

\cvitemit{2009 - presente}{Membro dell'Associazione ``Gruppo di Misure Elettriche e Elettroniche'' (\sbseries GMEE)}
\cvitemit{2023 - presente}{Membro della Commissione per la revisione del Sistema Informativo Integrato del  {(\sbseries GMEE)} - Associazione ``Gruppo di Misure Elettriche e Elettroniche''}
\cvitemen{2009 -}{Member of the Association ``Group of Eletrical and Electronic Measurements'' (\sbseries GMEE)}

\cvitemit{2009 - presente}{Membro dell'Institute of Electrical and Electronics Engineers (\sbseries IEEE)}
\cvitemen{2009 -}{Member of the Institute of Electrical and Electronics Engineers (\sbseries IEEE)}
\cvitemit{2011 - presente}{Membro del IEEE-IMS Technical Committee on Measurements and Networking, \sbseries IEEE IMS TC37}
\cvitemen{2011 - presente}{Member of IEEE-IMS Technical Committee on Measurements and Networking, \sbseries IEEE IMS TC37}
\cvitemit{2012 - presente}{Membro del IEEE-IES Technical Committee on Factory Automation, \sbseries IEEE TCFA}
\cvitemen{2012 - presente}{Member of the IEEE-IES Technical Committee on Factory Automation, \sbseries IEEE TCFA}
\cvitemit{2016 - presente}{Membro del IEEE-IES Technical Committee on Industrial Informatics, Sub-Committee Industrial Internet of Things, \sbseries IEEE TCII}
\cvitemen{2016 - presente}{Member of the IEEE-IES Technical Committee on Industrial Informatics, Sub-Committee Industrial Internet of Things, \sbseries IEEE TCII}

\cvitemit{2016 - 2018}{{\sbseries IEEE P61158} - IEEE Working Group P61158 per la definizione di uno standard di comunicazioni real-time industriali. Definitivamente pubblicato nel 2018 lo Standard IEEE 61158 Industrial Hard Real-Time Communication.}

\cvitemen{2016 - 2018}{{\sbseries IEEE P61158} - IEEE Working Group P61158 for the definition of a standard for industrial real-time communications. The IEEE 61158 Industrial Hard Real-Time Communication Standard was officially published in 2018.}

%\cvitemit{IEEE TCCC}{Membro del IEEE Computer Society Technical Committee on Computer Communications}
%\cvitemen{IEEE TCCC}{Member of the IEEE Computer Society Technical Committee on Computer Communications}

%\cvitemit{IEEE TCRTS}{Membro del IEEE Computer Society Technical Committee on Real-Time Systems}
%\cvitemen{IEEE TCRTS}{Member of the IEEE Computer Society Technical Committee on Real-Time Systems}

%\cvitemit{SPS Italia}{Membro del Comitato Scientifico di SPS Italia (https://www.spsitalia.it/it/comitato-scientifico)}
%\cvitemen{SPS Italy}{Member of the Technical Committee of SPS Italy  (https://www.spsitalia.it/it/comitato-scientifico)}
